% Generated draft: requires mathematical and source review.
\documentclass[11pt]{article}
\usepackage[margin=25mm]{geometry}
\usepackage{fontspec}
\setmainfont{Latin Modern Roman}
\usepackage{amsmath,amssymb,mathtools,unicode-math}
\setmathfont{Latin Modern Math}
\usepackage{enumitem,needspace,xurl,hyperref}
\setlength{\parindent}{0pt}
\setlength{\parskip}{5pt}
\setlength{\emergencystretch}{3em}
\newcommand{\N}{\mathbb{N}}
\newcommand{\Z}{\mathbb{Z}}
\newcommand{\Q}{\mathbb{Q}}
\newcommand{\R}{\mathbb{R}}
\newcommand{\C}{\mathbb{C}}
\newcommand{\F}{\mathbb{F}}
\newcommand{\A}{\mathbb{A}}
\newcommand{\PP}{\mathbb{P}}
\newcommand{\E}{\mathbb{E}}
\newcommand{\eps}{\varepsilon}
\newcommand{\dd}{\,\mathrm{d}}
\begin{document}
\section*{UnsolvedMath ID: 20000835}
\textbf{Source problem code:} catalogue entry
\subsection{Definitions and mathematical statement}
This record is a malformed aggregate, not one conjecture. The supplied verification says it begins in a citation at the end of source Problem (9), then concatenates Problems (10)--(19), with headers, control characters, and damaged displays. The phrase ``Conjecture 1'' belongs to reference [Mor05], not a new heading. The surviving questions are retained separately below under the existing permanent identity; no single truth claim resolves the aggregate.
\begin{enumerate}
\item The tail of (9) asks to prove Morita's conjecture, citing [Mor05], Conjecture 1, p. 4, and suggests a connection with Faber's intersection-number conjecture. Neither conjecture's full formulation is present, so the missing antecedent cannot be reconstructed reliably.
\item (10), Mondello: for a locus denoted \(\mathcal M_{g,n}^{\leq k\text{ rational components}}\), ask whether the indicated tautological degree \(R^{g-1+k}\) is \(\Q\), what generates it, when a one-dimensional socle is expected, and what the full tautological ring should be. The compactification, precise locus, and Chow-versus-cohomological grading are not recoverable from the damaged notation.
\item (11), Igusa: ask for operations relating stable classes on \(B\Gamma_\infty\), where \(\Gamma_\infty\) is the stable mapping class group, and for stable self-maps of the Thom spectrum \(\mathbb{CP}^{\infty}_{-1}\). The source writes
\[
H^*_{\mathrm{spec}}(\mathbb{CP}^{\infty}_{-1})\cong\Z[c_1]u,
\qquad |c_1|=2,\quad |u|=-2,
\]
and asks how this appears at the infinite-loop-space level. Exact stabilization and map conventions are not supplied.
\item (12), Baldwin: compute intersection cohomology of the moduli space of degree-\(d\) genus-\(g\), \(n\)-marked stable maps to \(\PP^r\), whose overline/compactification notation is damaged, and seek a suitable tautological ring. The virtual fundamental class is suggested as an analogue of the ordinary fundamental class, not asserted to give every required intersection-cohomology structure.
\item (13), Faber: determine the exact integral-cohomology order of the top Hodge Chern class \(\lambda_g\) on \(\mathcal A_g\), the moduli of principally polarized abelian varieties. The source attributes rational vanishing and a computation up to a factor of two to Ekedahl and van der Geer, and refers to Problem (5), absent here. Stack/coarse-space and coefficient conventions must be fixed before assigning an integral order.
\item (14), Bertram: a historical question about publication of Getzler's paper on \(\mathcal M_1\), with the playful ``\(t\to\infty\)'' comment and Getzler's objection to the phrasing. This is not a mathematical asymptotic conjecture to formalize as a theorem.
\item (15), Ellenberg: ask whether Hurwitz-space cohomology stabilizes as genus grows with covering degree fixed. Branching data, connected components, base curve, and stabilization maps are not specified. The motivation is degree-two point-count behavior and work on number fields/function fields in the 1980s attributed to ``Darskovksy and Wright'' in the extraction. More broadly, if normalized counts \(\#X_n(\F_q)/q^{\dim X_n}\) converge for all \(q\), ask whether a Harer-type stability mechanism explains them; no converse theorem from counts to stability is asserted.
\item (16), Tseng: ask a related stability question for maps \(C\to BG\) with finite \(G\). The source then separately reports Tillmann's connected-group comparison, for example \(G=S^1\), using
\[
E\operatorname{Diff}(F_{g,1})\times_{\operatorname{Diff}(F_{g,1})}
\operatorname{Map}_{\partial}(F_{g,1},BG).
\]
Here \(F_{g,1}\) is an oriented genus-\(g\) surface with one boundary component and boundary conditions are fixed. Gluing tori is said to induce homology isomorphisms in a range, but that range and exact hypotheses are absent. Finite and connected groups are different cases; the comparison does not settle the finite-group question.
\item (17), Sullivan: for a fixed compact connected complex curve \(C\), ask whether its connected finite unbranched covers become uniformly dense, in Teichmüller distance, in universally defined regions \(U_g\subset\mathcal M_g\). The literal quantifier is: for every \(\varepsilon>0\), for all sufficiently large \(g\), every point of \(U_g\) is within \(\varepsilon\) of a genus-\(g\) cover of \(C\). The source does not define \(U_g\). It connects this with finding arbitrarily close covers of two curves, called the Siegel--Ehrenpreis problem.
\item (18), Morita: compare the signature of a moduli-space cohomology pairing, the rational \(L\)-genus from Thom's rational Pontryagin classes, and the orbifold \(L\)-genus from orbifold tangent Chern/Pontryagin classes. It asks whether these agree, speculates that discrepancies reflect geometric versus complex-analytic singularities, and suggests analogous tautological-ring questions. The damaged compactification notation and omitted compactness/orientation conventions prevent definite numerical formulas; the proposed singularity interpretation is motivation, not a theorem established here.
\item (19), Sullivan: for \(LM=\operatorname{Map}(S^1,M)\), ask for the full algebraic structure relating a loop product and coproduct on homology, and whether working in spaces over \(M\) clarifies their Hopf/Frobenius aspects. The source refers to \(LM\times_M LM\), \(LM\times LM\), the evaluation map, and the diagonal of \(M\), but its diagram is corrupted. Orientability, coefficients, degree shifts, and local systems are not specified. The ordinary loop-fibration spectral sequence has \(E^2_{p,q}=H_p(M;H_q(\Omega M))\); a tensor-product description needs additional hypotheses. The extracted assignment of Hopf and Frobenius structures to base/fiber is not promoted to a general fact, and the missing diagram is not invented.
\end{enumerate}

The attached aid addresses only a quantifier obstruction in (17). If the base curve has genus \(h\), a connected unbranched cover of degree \(n\geq1\) has
\[
g-1=n(h-1).
\]
It reports the exact possible genus sets as \(\{0\}\) for \(h=0\), \(\{1\}\) for \(h=1\), and \(\{1+n(h-1):n\geq1\}\) for \(h\geq2\). Hence, if the regions \(U_g\) are eventually nonempty, the literal all-large-\(g\) density request is impossible unless \(h=2\). Genus two removes this arithmetic obstruction but does not prove density. A repaired question restricts \(g\) to the admissible spectrum.

For two hyperbolic bases of genera \(h_1,h_2\geq2\), the aid reports common admissible cover genera exactly when \(g>1\) and
\[
g\equiv1\pmod{\operatorname{lcm}(h_1-1,h_2-1)}.
\]
This is a genus-availability assertion, not existence of a common isomorphic cover or arbitrarily close covers. The partial obstruction does not resolve the remaining source questions, identify the regions, or supply current literature statuses. All aid results remain unreviewed supplied content.
\subsection{Short English statement}
The record concatenates the tail of one source question and ten others. Each surviving topic is retained with its attribution and missing conventions. The supplied aid only repairs a genus quantifier in Sullivan's cover-density question; it is not a solution of the aggregate or of its density assertion.
\subsection{Sources}
\begin{itemize}
\item AIM, \emph{A list of open problems and questions on the moduli space of curves}, workshop \emph{Topology and geometry of the moduli space of curves}: \url{https://aimath.org/WWN/modspacecurves/modspacecurves.pdf}. The supplied verification identifies the aggregate as the end of (9) followed by (10)--(19) in the five-page list.
\item Original attributions: Mondello, Igusa, Baldwin, Faber, Bertram, Ellenberg, Tseng, Sullivan, and Morita, with comments credited to Getzler and Tillmann. The tail cites [Mor05], Conjecture 1, p. 4; no full bibliography entry is supplied. Ekedahl--van der Geer and the extracted ``Darskovksy and Wright'' are mentioned without full citations.
\item Local registry \texttt{data/unsolvedmath/problems.json}, permanent ID \(20000835\), exact source code \texttt{catalogue entry}; canonical \texttt{aim-arithmetic-geometry-notes.json}, notes[80].
\item The attached research attempt supplies the cover-genus obstruction and no complete literature survey for the concatenated questions. It is explicitly machine-generated, not peer reviewed, and requires independent verification before citation.
\end{itemize}
\end{document}
